\chapter[Why Poetry Speaks Indirectly]{Why Poetry Does Not Say Things Straight}
\section{A Page Only Half Used}
Pick something up: an open poetry collection. Any collection will do, ancient or modern, Chinese or translated.

Do not read the words yet. Look at the paper. You will notice something almost absent from other books: expanses of blank space. A line often has only five or six Chinese characters, perhaps a dozen at most, before stopping mid-thought and starting another line. The words occupy a narrow strip, leaving white space beside them and between the lines.

A printer producing novels this way would distress the owner: does paper cost nothing? Write like this in an exercise book and a teacher's red pen will tell you not to waste paper. Novels, manuals, contracts, and news reports squeeze every page full. Only poetry confidently leaves blanks, stops in the middle of sentences, and refuses to finish saying things.

Compare a food-delivery app: "Spend thirty yuan, save five; delivery extra." Every word heads straight for its purpose. Both texts use the same language and Chinese characters, yet seem different species. One strives to be complete, precise, and quick; the other deliberately slows down, takes detours, and stops where stopping seems least appropriate.

Here is our question: \textbf{why does poetry not say things straight?} Are its blank spaces, line breaks, rhymes, and winding metaphors pretentious packaging, or another mode of speech---one that cannot speak in the language of "spend thirty, save five"?

\section{A Mutter During Morning Reading}
Come into a classroom with me.

Time: 7:25 on a December morning. Place: Class 3, second year, at a northern Chinese middle school. Daylight is incomplete, fluorescent tubes hum overhead, and condensation coats the windows. A pupil draws a smiling face in the mist, then wipes it away. The radiators are too hot to touch. Steamed buns scent the room.

The Chinese teacher leads morning reading. Today it is Li Bai's "Quiet Night Thoughts", a poem you may have memorised in nursery school:

\begin{quote}
Before my bed, bright moonlight; I wonder if frost lies on the ground. Raising my head, I gaze at the bright moon; lowering it, I think of home.
\end{quote}
More than forty voices recite in a drawn-out chorus, stretching guang, "light", and shuang, "frost", like a tail that will not flick free. On the third repetition, a mutter comes from the back, quiet but audible nearby. It is Big Liu, the class's strongest science pupil:

"If he misses home, why not just say so? Wo xiang jia le: 'I miss home', four characters, done. Those first three lines are pure detour."

Some laugh; some nod. The laughter is quiet but genuine. Perhaps you have thought something similar on an afternoon when failure to memorise a passage meant copying it as punishment.

The teacher hears. She neither gets angry nor lectures him about timeless masterpieces. Putting down the book, she says, "Fine. Let's try. Everyone rewrite the poem Big Liu's way: one sentence, all the meaning, not one unnecessary word."

After a brief silence, the class representative offers what everyone agrees is a perfect answer:

"Li Bai saw the moon at night and missed home."

The teacher writes it beside the original poem. "Here are two texts with exactly the same meaning. The one on the left has travelled thirteen hundred years. As for the one on the right, who would memorise it and tell their family tomorrow?"

No hands rise. Suddenly the steamed-bun smell seems very distinct.

Notice what happened. The teacher did not call Big Liu wrong. In \textbf{information}, his sentence scores full marks: time, person, event, feeling, nothing missing. Yet everyone senses that something has vanished, not a small detail but almost everything. Two texts with the same meaning sit together: one stirs you; the other resembles a parcel label.

That morning, the class was doing serious work: searching for what had disappeared. We will continue the search.

\section[The Plain-Speaking Challenge]{Take the Plain-Speaking Challenge Seriously}
Lay out the disagreement before choosing sides.

Big Liu's argument is stronger than it first appears and deserves its fullest statement. This practice of presenting an opponent's best case will serve us again: it prevents disagreement being dismissed as stupidity.

Plain speakers have at least three arguments.

First, \textbf{efficiency}. Language transmits information; tools should work efficiently. Wo xiang jia le, "I miss home", conveys in four characters what twenty characters and a moon convey, taking one-fifth the time. If detours are good, why not poetic weather forecasts? "Cloudy tomorrow with light rain, like heaven's half-confessed troubles." Try that and half the city will forget umbrellas.

Second, \textbf{fairness}. Poetry relies on hints, associations, and invitations to savour meaning, but savouring takes training. Those raised on poetry hear what lies between the lines; others remain outside. Obscuring important messages effectively charges admission, paid first by privileged households. Plain speech is communication's universal schooling: understandable regardless of background.

Third and sharpest, \textbf{protection against deception}. Ornate language has historically packaged more than it has expressed. Defeat becomes "strategic withdrawal", redundancies become "graduation", emptiness becomes "depth". Those with least to say most need pretty words. The more society tolerates indirectness, the more room deceivers gain. Plain speakers would rather inconvenience poetry than leave a back door for empty talk.

Each argument has everyday evidence. Dismissing them as science pupils' inability to appreciate beauty is dishonest.

Poetry's defenders have an equally complete case, also in three layers.

First, \textbf{some messages die when stated directly}. "I miss home", four Chinese characters, reports a fact but cannot transmit the feeling: waking at midnight to an oddly bright room, sleepily mistaking light for frost, reaching into cool moonlight, then realising that home is absent beneath this moon. Remove any link and the chain of feeling fails. Reporting only the conclusion hands over a cake's nutrition label: all the data, none of the cake.

Second, \textbf{slowness is a function, not a fault}. Delivery apps seek speed because their information matters for half an hour. Some words should remain in another person's heart longer. A eulogy does not say merely "He died; our condolences", nor a toast "Make money; cheers". At solemn moments humans slow down and attend to wording and rhythm. Taking time declares that something deserves the effort of saying.

Third, \textbf{a detour can be more honest}. Many genuine feelings are vague, uncertain, and contradictory. Which is more truthful: "I miss you", or "I stood looking at the moon for a strangely long time"? The first is tidy but clips a fuzzy thing into a square. The second winds around, yet its winding matches the feeling itself. Directness can distort such experience.

Now the real question emerges, not science versus humanities or useful versus useless:

\textbf{Are rhythm, rhyme, line breaks, metaphor, and silence, beyond a poem's literal information, wrapping paper or the goods themselves?}

If wrapping, Big Liu is right: say the thing and leave. If the goods, reducing the poem to one sentence destroys rather than compresses it. Voting cannot decide. We must open the poem and inspect its contents.

\section{Before Writing, Poetry Was a Memory Stick}
Before dismantling poetry, hear voices from different positions. Why poetry takes its form receives very different answers depending on where you stand.

First comes \textbf{the oral world}, a standpoint easily forgotten today. Poetry is far older than writing. Before paper, pens, or hard drives, genealogies, laws, remedies, and war histories lived in human heads. Heads are unreliable: loose sentences slip away or become muddled. Civilisations independently invented the same technique: \textbf{turn important words into rhythmic, rhyming, patterned songs}.

Before being written, the Iliad and Odyssey were sung for centuries. Their many thousands of lines relied on formulas, repeated ready-made phrases such as "wine-dark sea" and "rosy-fingered dawn". With words ready on the tongue, singers could plan the next passage. At the Iliad's opening, the singer calls on a goddess to sing Achilles' anger. This is more than courtesy: a start-up ritual tells listeners that memory mode is running.

On China's Tibetan Plateau, performers still recite the Epic of King Gesar for weeks, its scale exceeding any fixed written epic, much of it living only in voices. West African griots likewise preserve royal genealogies and histories across generations. They are walking libraries, and rhyme, rhythm, and formula are their shelves.

If lessons resist memorisation while lyrics stick after two hearings, it need not mean laziness. An ancient machine inside you still works: rhythm and melody make words memorable. \textbf{Many poetic oddities---rhyme, repetition, regular beat---began as storage technology, not aesthetics.} That explains their existence, though not yet their emotional power. One layer at a time.

The second voice comes from \textbf{court and temple}. The Book of Songs divides into feng, "Airs", ya, "Odes", and song, "Hymns". Hymns addressed ancestors and deities; many Odes accompanied court banquets. Rulers used poetry not principally to move hearts but to establish order: which songs belonged to which occasions, and who could stand where. Indirectness became power's language: solemn, elaborate speech deliberately distanced itself from market bargaining, announcing that words here were different.

The third voice comes from \textbf{fields and streets}. The Airs gather regional songs of lovers missed, forced labour resented, and faithless men rebuked. "Cutting Sandalwood" from the Airs of Wei confronts people with full granaries who do no work: why should those who neither sow nor reap take our grain? A workers' song over two thousand years old, with rhythmic anger. At court, indirectness established rules; in the fields, it offered camouflage. Insulting a lord directly might cost your head; singing retained the defence that it was only a song.

Hear the three answers? Singers, officials, and farmers speak indirectly to remember, command respect, or survive. None has the sole correct explanation. Together they refute poetry's supposed monopoly by refined scholars. Poetry's source code was written in throats and fields.

\section[Thinking Bridge: Poetry into Prose]{A Thinking Bridge: Compress Poetry into Prose and See What Is Lost}
Here is a portable tool requiring no special gift, only paper. It extends the morning teacher's method.

Call it the \textbf{compression test}: rewrite each poetic line as complete, fluent prose without losing meaning, then list what disappears in rewriting. That loss inventory reveals the poem's toolkit: what disappears was doing work.

Try Wang Wei's "Birdsong Brook", only twenty Chinese characters:

\begin{quote}
At ease, a person; cassia blossoms fall. The night is still, the spring mountain empty. Moonrise startles mountain birds; at intervals they call within the spring ravine.
\end{quote}
A full-information prose version: "Someone is at leisure while cassia blossoms fall. The night is quiet and the spring mountain empty. The moon rises, disturbing mountain birds, which occasionally call beside the stream."

All the meaning remains. Now list the losses.

\textbf{Loss 1: Sound.} The Chinese sounds xian, "at leisure", shan, "mountain", and jian, "between", quietly echo one another. Luo, "fall", chu, "emerge", and yue, "moon", are short checked-tone syllables, pebbles dropped into water, while most sounds are light and level. The twenty characters form a soundscape of prevailing stillness and occasional movement. The prose sounds like ordinary speech; that soundscape disappears.

\textbf{Loss 2: Rhythm.} Each original line has five characters, making breath fall into four equal rhythmic sections, like a leisurely walker: four steps, pause, four steps, pause. This unhurried pace physically embodies leisure and stillness. You do not merely read quiet; you \textbf{experience it through your mouth}. Prose breaths vary in length, turning strolling into hurrying.

\textbf{Loss 3: Lines and pauses.} Ren xian gui hua luo, "At ease, a person; cassia blossoms fall", stands alone in five characters. The pause lets the falling blossoms briefly illuminate your mind without the next image covering them. Prose makes them a passing clause, gone before they shine.

\textbf{Loss 4: Word order and omission.} The poem names no subject: who sees the blossoms? Who is the person at leisure, poet or reader? It does not say, leaving a seat anyone may occupy. Prose dutifully supplies the subject, making the sentence fluent but welding that seat shut. More strikingly, the poem never says "quiet" outright. It offers blossoms, empty mountains, moon, and a birdcall, using sound to bring out silence. Letting the conclusion grow inside the reader is \textbf{leaving space}.

Four losses reveal four tools. Give each a label and an everyday example:

\begin{itemize}
\item \textbf{Rhythm}: Language's walking pace and heartbeat. You speak quickly when angry and slowly when soothing a child. Poetry makes an existing practice into craft.
\item \textbf{Prosody, including rhyme and tone patterns}: Echoes between sounds. Hear the Chinese advertising line "No gifts this holiday" and the next line arrives automatically: sound hooks words into memory.
\item \textbf{Lineation and pauses}: Control which image you see and for how long, like film editing: each cut creates a shot.
\item \textbf{Leaving space}: Deliberate omission for you to complete, like a storyteller pausing at a joke's crucial moment. You finish the joke when you laugh.
\end{itemize}
Five frequent companions join these tools: \textbf{imagery}, concrete pictures carrying feeling, as the moon carries Chinese memories of home beyond astronomy; \textbf{metaphor}, directly calling one thing another, as "Time is a thief" is bolder than "Time is like a thief"; \textbf{symbolism}, a concrete thing conventionally representing an abstraction, like a dove for peace; and \textbf{repetition}, returning words gaining weight like a refrain.

The compression test turns poetic quality from mystery into engineering. Instead of arguing about atmosphere, compress and inspect the losses. A longer list suggests more components working in the original.

Here is an \textbf{easily misread counterexample}. You might conclude that poetry dresses ordinary sentences in pretty words and chops them up. Try breaking medicine instructions into short lines: "Three times daily / take with warm water / avoid spicy food". It looks poetic but is not poetry. Conversely, William Carlos Williams's "The Red Wheelbarrow" uses no fancy words: many things depend on a red wheelbarrow wet with rain beside white chickens. It is recognised as poetry because its line breaks make you inspect the wheelbarrow piece by piece and see an overlooked brightness. Together these tests show that \textbf{poetry lies neither in ornament nor layout, but in whether its tools actually work}. Pretty replacements alone usually produce a parcel label decorated with icing.

\section[Three Traditions, Three Detours]{Three Traditions, Three Kinds of Indirectness}
Carry the toolkit through several poetic traditions. The same parts become different machines, warning us against treating any one tradition as poetry's standard form.

\textbf{First stop: classical Chinese poetry builds with sound.} Chinese has tones. Regulated patterns prescribe alternation between level tones, sustained and even, and oblique tones, clipped or turning: \textbf{tonal prosody}. Add \textbf{parallelism}, matching word classes and mirrored structures across paired lines. Du Fu's "Quatrain" offers "Two yellow orioles sing in emerald willows; a line of white egrets rises into blue sky." Two matches one line, yellow orioles white egrets, sing rise, emerald willows blue sky. The fit is exact without stiffening the image. Such form produces \textbf{poetic atmosphere}: a few objects, supported by sonic symmetry and visual parallels, open a mental space larger than their literal meaning. Twenty Chinese characters can contain a mountain partly by leaving things unsaid.

\textbf{Second stop: European sonnets chain rhymes together.} Fourteen lines, roughly ten syllables each, with prescribed rhyme positions. Shakespeare wrote 154; a famous opening asks:

\begin{quote}
Shall I compare thee to a summer's day?
\end{quote}
It asks whether the beloved resembles summer, then answers no: the beloved is lovelier, because summer passes while the poem preserves its subject. Notice the structure: fourteen lines make a small room where a comparison is built, dismantled, and transformed by the final couplet. Form is not merely a container; it determines where thought turns.

\textbf{Third stop: Japanese haiku sharpens the pause.} Haiku is the world's shortest fixed poetic form, seventeen sound units arranged five, seven, five, with a kigo, a seasonal word. The most famous, by seventeenth-century Matsuo Basho, reads:

\begin{quote}
An old pond---a frog jumps in---the sound of water.
\end{quote}
Its sense is simply an old pond, a leaping frog, and water sounding. Almost nothing is stated before it ends. Yet the break after the first phrase, the Japanese kire or cut, like a freeze-frame, divides worlds before and after the splash. Before, a thousand years of stillness; afterwards stillness remains, but you have changed. Haiku pushes omission nearly to stinginess.

\textbf{Fourth stop: modern free verse dismantles rules but keeps lines.} Nineteenth-century American poet Whitman declares at the opening of Leaves of Grass:

\begin{quote}
I celebrate myself, and sing myself.
\end{quote}
He celebrates and sings his own being. Unrhymed, uneven sentences resemble prose: free verse removes inherited rhyme and metre. Yet removing the frame does not discard the tools. Whitman's \textbf{repetition} and parallel clauses arrive like waves, while lineation controls the eye. Free verse invents a temporary form for each poem rather than having none. Its apparent freedom requires decisions at every step, demanding even more craft.

After four stops, the plain-speaking challenge looks different. Rhythm, rhyme, lines, and prosody are not pretty substitutions for ordinary sentences. They work simultaneously through sound, page, and meaning: sound addresses body and memory, layout guides eyes and pauses, meaning invites association and omission. A good poem draws all three tight at once.

\section[Evidence: An Ancient Love Song]{A Little Evidence: A Love Song Twenty-Five Centuries Old}
Theory cannot replace close reading. This poem from the Zhou Nan section of the Book of Songs belongs to China's earliest surviving poetry, roughly twenty-five to thirty centuries old:

\begin{quote}
Guan-guan cry the ospreys on an island in the river. The gentle, lovely maiden is a fine match for the gentleman.
\end{quote}
("Guan Ju", Zhou Nan, Book of Songs.)

The meaning: ospreys call guan-guan on a river islet; the quiet, admirable girl makes a worthy partner for a gentleman.

Apply the compression test: "A gentleman wants to marry a good girl." Seven Chinese characters, with no meaning lost. Now the losses:

\textbf{Sound.} Read the original's sounds aloud. Guan-guan repeats a birdcall; ju-jiu, "osprey", contains related sounds, as does hao-qiu, "fine mate", an echo called repeated rhyme. The opening is a sound game, with cries and similar syllables colliding and answering. A springtime birdcall from twenty-five centuries ago remains stored in these six characters, sounding again when read today. No paraphrase captures it.

\textbf{An odd beginning.} A love song begins with birds. Ancient critics called this \textbf{xing, an evocative opening}: sing about something else before approaching the subject. A blunt proposal risks abrupt rejection. Begin instead with birds, creating springtime, riverside, paired companionship; let listeners enter that feeling before people appear. The osprey is no random choice: tradition describes faithful mating pairs. The bird is not just introductory music but the first metaphor.

\textbf{A vacant seat for anyone.} There are no names, no named location beyond a river, no backstory. Who are the gentleman and maiden? We do not know. That allows readers across twenty-five centuries to put their stories into these sixteen characters. Omission makes room: the less specified, the more people can enter.

A candid aside: generations of interpretation have wrapped this poem in layers. Han scholars insisted it praised a consort's virtue, a moral lesson for royal women; after the Song, others restored it as a folk love song. The same sixteen characters gave the court rules and ordinary people heartbeats. The disagreement teaches that \textbf{meaning lives not only in words but in the reader's position}. We explore that next.

\section[One Poem, Three Explanations]{One "Quiet Night Thoughts", Three Explanations}
Return to the classroom poem. With tools available, compare genuinely different explanations of its power. Separate four things: what happened, a twenty-character poem surviving thirteen centuries and memorised by hundreds of millions, is factual; why it moves us is interpretive and disputed; how contemporaries understood it is another question, and Li Bai cannot answer us; whether it deserves masterpiece status is a value judgement. We compare the second layer.

\textbf{Explanation 1: Content captures a universal feeling.}

The poem moves us by expressing something everyone feels but cannot always articulate: at night, moonlight enters and homesickness follows. Everyone shares the moon; everyone is vulnerable to missing home. The poem hits the centre. This is simple and intuitive: remembering it during your first homesickness offers evidence. But thousands of poems concern home, and countless more the moon. Why did these twenty characters win? If subject matter were everything, all equally accurate hits would be equally famous. Content alone cannot explain this centuries-long contest.

\textbf{Explanation 2: Sound and structure generate the power.}

Use our toolkit. Say guang, "light", shuang, "frost", wang, "gaze", and xiang, "home": their broad shared rhyme opens the mouth and carries far, bright and spreading like moonlight. The first two lines create a cool illusion, mistaking moonlight for frost; the last pair two movements, raising and lowering the head. Your neck travels from sky to home. No word says cold or loneliness, yet the illusion's chill and lowered head say everything. This explains why one poem wins among similar subjects: every formal part works. Its limit is explaining why it works this extraordinarily well. Copy the blueprint into a rhyming five-character quatrain and mediocrity remains possible. Plans cannot guarantee a masterpiece.

\textbf{Explanation 3: Cultural memory supplies two millennia behind the moon.}

Zoom out. Your response to "Before my bed, bright moonlight" does not arise between reader and poem alone; a whole culture intervenes. You have watched Mid-Autumn moons, seen mooncake-box moons, heard "The moon is brighter back at home", and memorised "A bright moon rises over the sea". Poems, songs, festivals, and reunion meals charge the image until its appearance lights an entire memory network. Li Bai did not invent the circuit; he joined its largest branch. This explains something awkward for the other accounts: someone outside Chinese culture may understand every translated word yet feel little, lacking not vocabulary but that circuit. Its weakness is circularity: famous because everyone learns it, learned because famous. It also unfairly makes poetry a parasite fed by tradition. Someone had to lay the first wires.

None settles the bill alone. Content explains what is said; form how; cultural memory why it strikes you in particular. This is not fence-sitting but a methodological lesson: \textbf{when an explanation claims everything, ask what it left out}.

\section[Further Challenge: Defamiliarisation]{A Deeper Challenge: Defamiliarisation Makes Stone Stony}
Now bring in our strongest theory. It addresses a question we have avoided: why is compression into prose fatal to poetry?

Early twentieth-century Russian literary scholars later called Formalists included Viktor Shklovsky. His combative 1917 essay "Art as Technique" struck a conceptual match still burning more than a century later.

He began with an everyday observation: \textbf{perception grows rusty}. Learning to cycle demands complete attention; later you ride a street while thinking of dinner. How many stairs lead to your flat? Despite daily use, you probably cannot say. Habit automates perception. After ten thousand encounters, you see not the thing but a label. Most of the time, Shklovsky argues, we recognise rather than see the world, then move on.

Art's task is \textbf{to resist this automation}. Its method, ostranenie, translated as defamiliarisation or making strange, presents familiar, supposedly known things through unfamiliar difficulty and detour, forcing another look: real sight rather than recognition. His famous formulation says, roughly, that art restores our feeling for things and makes stone more stony.

Now poetry makes sense. Why not speak plainly? \textbf{Because plain speech is the escalator you have ridden ten thousand times: you arrive without moving or seeing anything.} Poetry removes it and supplies unfamiliar stairs requiring lifted feet: an awkward metaphor, an unexpected break, or Du Fu's yong, "surge", in "The moon surges as the great river flows", from "Thoughts While Travelling at Night". How can moonlight surge like water? That second of hesitation makes it visible. Artistic language, for Shklovsky, is obstructed, distorted, prolonged, delaying arrival at meaning because feeling happens along the journey, not at the destination.

Why does "Li Bai saw the moon at night and missed home" leave nothing? It gives the destination's address and removes the road.

This powerful theory also explains unusual camera angles, arresting song introductions, and jokes more memorable than press releases. But it faces two challenges you should learn to ask.

\textbf{Challenge 1: Defamiliarisation itself becomes automatic.} A technique opens eyes once and becomes routine after ten thousand uses. Martial-arts novels' "dark moon, high wind" once startled; now it may amuse. Internet slang first estranges familiar phrases, then after six months becomes more automatic than "very good". Defamiliarisation is no permanent recipe but an endless arms race: language wears out and poets need fresh ammunition.

\textbf{Challenge 2: Can it excuse incomprehensibility?} If more obstruction means more art, would torn dictionary scraps randomly glued together make the greatest poem? Obviously not. The distinction is \textbf{whether difficulty pays off}. Good estrangement leads through a detour to an unseen landscape; bad estrangement ends nowhere, or somewhere even its author cannot explain. Difficulty is cost, the view is return. This accounting is sounder than equating clarity with shallowness and obscurity with depth. Writing shielded by "You lack the ability to understand" rarely survives that audit.

The theory leaves Big Liu an honourable role. He is no enemy of poetry, but a sensitive detector of \textbf{failed defamiliarisation}. His alarm is right when poems stack worn phrases like falling blossoms and broken hearts. Those are not poetic stairs, but automatic escalators disguised as them.

\section[Headphones, Ads, and Chats]{In Headphones, Advertisements, and Chat Histories}
This ancient question is busy throughout daily life, whether you notice or not.

\textbf{Start with your headphones.} You may struggle to memorise "The Yueyang Tower", yet sing dozens of songs perfectly. That ancient memory machine works for you: rhythm, rhyme, and melody weld lyrics into the mind. Rap takes old rhyming crafts into radical double and triple rhymes; flow, words' placement against beats, is rhythm's frontier laboratory. In a good verse, stressed syllables, breaks, and sudden accelerations use our tools in new clothes. Next time rap gives you goosebumps, compress it into smooth prose and read it aloud. How many goosebumps survive?

\textbf{Look at the street.} Advertising spends lavishly studying indirect speech. Good copy does not merely say a product offers quality and value; that is a parcel label. It gives imagery, rhythm, and a hook, letting the conclusion grow inside you. Repeated promotional broadcasts use \textbf{repetition}; beach-and-drink advertisements use \textbf{associations between images}. Poetic craft has not disappeared; it changed jobs and earns well. This answers plain speakers' third concern: yes, indirectness packages products. All the more reason to learn its tools and resist its hooks.

\textbf{Then comes translation, an unavoidable scene.} Every foreign poem you read may be the remains of an impossible mission. Poetry works through sound, page, and meaning simultaneously; translation can almost never preserve all three.

Take Basho's haiku: furuike ya / kawazu tobikomu / mizu no oto, seventeen sound units in five-seven-five with a central cut. The translator faces cruel choices. Preserve that count in Chinese and words must be cut or padded, risking altered imagery. Preserve the old pond, frog, splash, and pause, and the sound-count structure slips away. Add pleasing Chinese rhyme, though the original does not rhyme, and is that translation or rewriting? Chinese versions range from tidy five-character lines to loose prose, with no standard answer. \textbf{Every translation is an open ledger recording what its translator chose to sacrifice}.

Shakespeare's comparison to a summer's day offers another problem. In England, summer is mild, bright, brief, and precious: supreme praise. A reader who has endured August in Wuhan may hear an insult. One image changes meaning with climate. Translate literally and preserve words while losing the English summer; replace summer with lovely Chinese spring and preserve effect while losing the original. This is no accidental error but a structural dilemma: \textbf{images carry local registration papers; many feelings cannot transfer residence}. Italians put it wittily: translator, traitor. Once you understand that, a celebrated translation prompts another question: what was betrayed, and what preserved?

\textbf{Finally, the newest variable: machines.} Today's AI produces batches of regulated Chinese verse with impeccable tones and rhymes and every expected image, falling blossoms, lone boats, slanting sunlight. Many readers respond: polished, pretty, then nothing. Audit it with our tools. Sound and page may score full marks, but Shklovsky's cost-return account often fails: those images are stairs worn flat by millions of repetitions, with nothing newly seen. The cultural-memory circuit connects statistically common plugs while bypassing your particular socket. That judgement may hold today without holding forever. Machines change, and force a once-unnecessary question: if metre, imagery, and atmosphere are computable, what uncomputed remainder makes poetry?

Another version quietly surrounds you. Young people reinvent lineation in messages: splitting one thought into seven or eight sends, breaking sentences, isolating important words. Nobody teaches this layout, yet everyone reads its pauses. Poetry's blank page spaces live again inside our typing boxes.

\section{Your Turn: Who Pays for the Detour?}
Return to the classroom: "Quiet Night Thoughts" on the left of the board, the homesickness summary on the right.

Now a new scene, and our final question.

Big Liu grows up and becomes a father. Away on business late one night, he sees moonlight through hotel curtains and misses home. Wanting to message his three-year-old daughter, he types "Daddy misses you", then deletes it. He types "The moon is especially bright here; can you see it there?", thinks, deletes again, and finally sends only a moon photograph, without words.

His daughter receives it next morning. \textbf{How old must she be to understand the photograph? Until then, has this wordless message said something or nothing?}

Before answering, weigh our arguments again.

Big Liu has a point: indirectness costs something, often paid by listeners. A three-year-old cannot afford the tuition of reading moons. For years the photograph may be ordinary to her. Clarity keeps the cost with the speaker; detours sometimes transfer it to others. The plain-speaking fairness argument still holds.

The other side has an account too. Some things shrink when spoken. "Daddy misses you" is true but cannot contain the man at the curtain: his hesitation, two rounds of typing and deletion, his final wordlessness. The photograph leaves space. He pays now by betting that one day, older and away from home beneath a moon, his daughter will suddenly be struck by it. She will then receive a parcel no direct statement could deliver. Poetry's accounting across time holds too.

The question is therefore harder than whether plain speech or detours are better: \textbf{may someone say something important in a way another person cannot yet understand?} Who pays, when, and on what grounds may we gamble that the listener can eventually afford it?

Give your answer and reasons. Notice yourself arranging language, choosing what comes first and where to pause: deciding rhythm, lineation, and omission. Whatever your conclusion, you are already answering with poetry's tools.
